\section{Conclusions}

This study evaluated the contribution of segmentation and explicit texture descriptors in lightweight CNN-based rice leaf disease classification through a controlled ablation design. The results show that the proposed segmentation-texture hybrid approach is feasible and competitive, but it was not the strongest configuration in the evaluated setting. Across five random seeds, the CNN with explicit GLCM+DWT texture fusion (E2) achieved the highest mean macro-F1 score, while the segmentation-only model (E1) achieved the lowest inference time. The proposed hybrid model (E3) improved over the baseline in aggregate accuracy, but it did not produce a statistically significant advantage over the baseline or outperform the texture-fusion configuration.

The main empirical finding is that, under a pHash-controlled split, explicit texture features provided the most consistent performance gain, whereas segmentation mainly contributed to efficiency. This suggests that handcrafted texture and frequency descriptors can still complement CNN embeddings in rice leaf disease classification, especially when symptoms are visually expressed through lesions, spots, and local texture patterns. However, the high baseline performance indicates that the dataset presents a performance saturation scenario, where additional components may not produce proportional or statistically significant improvements.

Therefore, the contribution of this work lies not in claiming absolute superiority of the proposed hybrid model, but in providing a controlled comparison of four configurations under the same split, training setup, metrics, and seeds. This ablation evidence helps clarify the relative effect of segmentation, texture fusion, and their combination in a lightweight CNN pipeline.

Future work should validate the same experimental design on external and field-acquired datasets with more diverse lighting conditions, backgrounds, cameras, and disease stages. Additional work should also explore stronger near-duplicate grouping strategies, explainability methods such as Grad-CAM, and real-time deployment evaluation on resource-constrained devices. These extensions would help determine whether segmentation and texture fusion provide larger benefits under stronger domain shift and practical agricultural conditions.
